
\subsection{End of Life}
\subsubsection{Metrics}
\begin{itemize}
\item Mass\\
\\
Mass is a critical component of any spacecraft. STARS must conform to specific mass requirements in order to match the 3U standard as stated in \textbf{FR.1} The mass of the CubeSat depends on the mass of each individual system. So, each system should aim to be as light as possible while ensuring that they meet all design requirements.
\\
\item{Volume} \\
\\
The entire size of the CubeSat shall be constrained to a 3U standard as specified by requirements \textbf{FR.1}  and  \textbf{DR.1.1.}  As such, a maximum of 1U or less has been allocated to the end of life systems on this mission. Any divergence from this strict standard would negatively affect ride share options and could possibly lead to the inability to launch.
\\
\item{Power Consumption} \\
\\
Electrical power is a limited resource on board any spacecraft. The STARS on board power system will only be able to supply a given amount of power at any given time. Ensuring that the EoL system requires as little power as possible reduces the strain on the overall power budget.
\\
\item{Cost} \\
\\
Budget constraints are an important part of system design. For the project to remain within budget, the propulsion system should strive to be as inexpensive as possible while still maintaining capability to perform all the desired tasks.  
\end{itemize}

Table \ref{tbl:EoL_Metric_Scores} below provides a description for how different values would be assigned to each design option in regards to each metric. 

\begin{table}[H]
\centering
\begin{tabular}{|l|c|c|c|c|c|} 
\hline
    \textbf{Metric} & \textbf{5} & \textbf{4} & \textbf{3} & \textbf{2} & \textbf{1}
    \\
\hline
    \textbf{Mass} & <100g & 100-150g & 150-200g & 200-300g & >300g 
    \\
\hline
    \textbf{Volume} & <75cm\textsuperscript{3} & 75-100cm\textsuperscript{3} & 100-150cm\textsuperscript{3} & 150-300cm\textsuperscript{3} & >300cm\textsuperscript{3} 
    \\
\hline
    \textbf{Power} & <1W & 1-5W & 5-10W & 10-25W & >25W 
    \\
\hline
    \textbf{Cost} & <\$10,000 & \$10,000-50,000 & \$50,000-\$100,000 & \$100,000-\$200,000 & >\$200,000 
    \\

\hline
\end{tabular}
\caption{EoL Trade Study Metric Values}
\label{tbl:EoL_Metric_Scores}
\end{table}
\subsubsection{Weights}
\begin{itemize}
    \item Mass - 0.35
    \item Volume - 0.35
    \item Power - 0.1
    \item Cost - 0.2
\end{itemize}

The Mass and Volume metrics are both given an equal weighting. This is because until activation, the EoL system is merely dead weight to the CubeSat. It is of paramount importance that the EoL system minimize its burden on the other CubeSat Systems throughout the mission phase. This way, STARS can more easily achieve its goals. For this reason, Mass and Volume are both rated equally, and they are both rated higher than the other two metrics. 

The Power metric is given only a 10\% weight, much smaller than the other three. This is because the EoL system will only be activated at the end of the mission, remaining in a dormant state until almost all of the other systems are likely to be inactive. At this time, the CubeSat will likely have more power available to send to the EoL system as, for example, the payload will no longer require power. Additionally, for all three of the passive options, drag chutes, electromagnetic tethers, and deployable booms, power is only required during deployment. This means that the power draw they require would only be for a very limited amount of time, likely on the magnitude of a few seconds to a few minutes. 

The Cost metric, while important, is not as critical as Mass and Volume, and as such is given a slightly lower weighting, at 20\%. Cost is, of course, an important thing to consider, the entire CubeSat must fall within a certain budget and minimizing the cost of any individual component or system helps ensure that this is the case. However, in the case of EoL systems, there is another aspect to consider. If the EoL system fails to deorbit the CubeSat within the timeline specified by NASA STD 8719.14B, there is a possibility that NASA or the IADC could impose strict fines on the organization responsible. Because of this, spending extra money to ensure that the EoL system is reliable and will preform adequately is reasonable. 
\subsubsection{Potential Pugh Matrix}

Table \ref{tbl:EoL_Pugh} below is an example Pugh matrix that could be generated using the metrics and weights defined above once significant analysis has been performed on each of the different Key Design Options. 
\begin{table}[H]
\centering
\begin{tabular}{|l|c|c|c|c|c|} 
\hline
    \textbf{Category} & \textbf{Weight} & \textbf{Cold Gas} & \textbf{Drag Chute} & \textbf{Electromagnetic Tether} & \textbf{Deployable Boom}
    \\
\hline
    \textbf{Mass} & 0.35 & 1 & 3 & 5 & 5
    \\
\hline
    \textbf{Volume} & 0.35 & 2 & 3 & 5 & 3 
    \\
\hline
    \textbf{Power} & 0.1 & 3 & 4 & 5 & 5 
    \\
\hline
    \textbf{Cost} & 0.2 & 4 & 4 & 3 & 4 
    \\
\hline
    \textbf{Total} & 1 & 2.15 & 3.3 & 4.6 & 4.1 
    \\
\hline
\end{tabular}
\caption{Example EoL Pugh Matrix }
\label{tbl:EoL_Pugh}
\end{table}

In this example case, the electromagnetic tether option has emerged as the best option, narrowly beating out deployable boom. This is a good indication that, should no other critical errors be uncovered with this kind of system, the electromagnetic tether is a good option for the EoL system.
